\section{Variational Autoencoder} \label{sec:vae}
A Variational Autoencoder (VAE) is a type of generative ML model based on a neural network architecture capable of learning an efficient and probabilistic encoding of data into a latent space of low dimensionality. By sampling vectors from the learned distribution in the latent space, a decoder can afterwards generate new and realistic data samples, even when the network was trained to encode complex and high-dimensional input data ~\citep{kingma2013}. The scheme in Fig.~\ref{fig:vae} shows a typical architecture of a VAE composed of convolutional layers. First, the encoder maps input vectors $\vec{x}$ to the parameters of a multi-dimensional probability distribution over the latent space, \eg a Gaussian with parameters $\vec{\mu}$ and $\vec{\sigma}$ (note that bold symbols denote vector quantities). Second, a vector $\vec{z}$ is sampled in the latent space and passed to the decoder, which aims to reconstruct the original input $\vec{x}$. 

\begin{figure*}[htbp]
  \centering
  \includegraphics[width=\textwidth]{f01}
  \caption{
    The basic scheme of a VAE composed of two-dimensional convolutional layers. 
    The encoder compresses an input $\vec{x}$ into the low-dimensional latent space.
    The decoder attempts to reconstruct the original data from a vector $\vec{z}$ sampled from the latent space.
  }
  \label{fig:vae}
\end{figure*}

The training of the VAE is guided by a loss function that balances the fidelity of the data reconstruction ($\mathcal{L}_R$)  against a term that regularizes the probability distributions in the latent space towards a simple prior distribution $p(\vec{z})$ ($\mathcal{L}_{KL}$):

\begin{equation}
  \mathcal{L}_{VAE} = -\mathbb{E}_{q_{\phi}(\vec{z}\vert\vec{x})}
  \left[ \log p_{\theta}(\vec{x} \vert \vec{z}) \right] + 
  D_{KL} \left( q_{\phi}(\vec{z}\vert\vec{x}) \Vert p(\vec{z}) \right)
\label{eq:L_vae}
\end{equation}
where $q_{\phi}(\vec{z}\vert\vec{x})$ is the encoder approximate posterior and $p_{\theta}(\vec{x} \vert \vec{z})$ the decoder likelihood. Equivalently, it can be expressed as the sum of two terms:
\begin{equation*}
    \mathcal{L}_{VAE} = \mathcal{L}_R + \mathcal{L}_{KL}
\end{equation*}
The first term ($\mathcal{L}_R$, reconstruction loss) encourages reconstructions to be accurate and, in this work, it is assumed to be the Mean Squared Error (MSE) between the input ($\vec{x}$) and its reconstructed output (\vec{x}'):
\begin{equation}
  \mathcal{L}_R = \dfrac{1}{\spaceDim} \sum_{i=1}^{\spaceDim} (x_i - x'_i)^2
\end{equation}
being \spaceDim the dimension of the original space. The second term ($\mathcal{L}_{KL}$) is defined in terms of the Kullback-Leibler (KL) divergence, $D_{KL}$, and measures the distance between the approximate posterior $q_{\phi}(\vec{z} \vert \vec{x})$ and the prior distribution $p(\vec{z})$. In practice, the $D_{KL}$ term forces $q_{\phi}$ to be close to the prior distribution $p(\vec{z})$. Traditionally, the standard VAE assumes a diagonal multivariate normal distribution for the approximate posterior, $q_{\phi}(\vec{z} \vert \vec{x}) = \mathcal {N}(\vec{\mu}, \vec{\sigma^2} \mathbb{I} )$, and a standard normal distribution for the prior, $p(\vec{z}) = \mathcal {N}(\vec{0},\mathbb{I})$, leading to:
\begin{equation}
  \mathcal{L}_{KL} = -\dfrac{1}{2} \sum_{i=1}^\latentDim \left[1+\log \sigma_i^2 - \mu_i^2 - \sigma_i^2 \right]
\end{equation}
being $\vec{\mu} = (\mu_1,\dots,\mu_\latentDim)^{\intercal}$ the vector of means, $\vec{\sigma^2} = (\sigma_1^2,\dots,\sigma_\latentDim^2)$ the vector of variances and \latentDim the dimension of the latent space where vectors $\vec{z}$ live. Finally, the so-called $\beta$-VAE is a modification of the standard VAE that uses a hyperparameter $\beta$ to weight the KL divergence term in the loss function:
\begin{equation}
  \mathcal{L}_{VAE} = \mathcal{L}_R + \beta \, \mathcal{L}_{KL}
  \label{eq:vae_loss}
\end{equation}
This is the expression of the loss function used in this work. In summary, VAE training optimises the encoder and decoder by minimising the VAE loss function \eqref{eq:vae_loss} on a training set of input vectors $\left\{ \vec{x_i} \right\}$, so that inputs are accurately reconstructed from latent samples, while the latent distribution remains close to the prior.